\documentclass[11pt]{article}

\usepackage[margin=1in, a4paper]{geometry}
\setcounter{secnumdepth}{0} % keine Nummerierung
\usepackage[T1]{fontenc} 
\usepackage[utf8]{inputenc}
\usepackage[ngerman]{babel}

% ==================== NATIVE SCHRIFTKOMBINATION ====================
% 1. Libertine für Fließtext (Serif)
\usepackage{libertine} 

% 2. Source Sans für Überschriften
\usepackage{sourcesanspro} 
% Alternativer Font: 
% \usepackage{cabin} -> % entfernen und % vor \usepackage{sourcesanspro} setzen

% 3. Semi-Bold:
\DeclareRobustCommand{\sbseries}{\fontseries{sb}\selectfont}
\DeclareTextFontCommand{\textsb}{\sbseries}
% ===================================================================

\usepackage{titlesec}

% Überschriften mit nativer Semibold-Sans
\titleformat{\section}
  {\large\sffamily\sbseries\uppercase}{}{0pt}{}[\vspace{1.3pt}\titlerule]

\titleformat{\subsection}
  {\large\sffamily\sbseries}{}{0pt}{}

\titleformat{\subsubsection}
  {\normalsize\sffamily\sbseries\itshape}{}{0pt}{}

% ==================== ABSTÄNDE (global) ====================
% Abstand zwischen Absätzen/Zeilen (kein manuelles \vspace{} zwischen den Absätzen)
\setlength{\parskip}{0.2em}

% Abstände der Überschriften: {links}{vor Überschrift}{nach Überschrift}
\titlespacing{\section}{0pt}{15pt}{10pt}
\titlespacing{\subsection}{0pt}{6pt}{4pt}
\titlespacing{\subsubsection}{0pt}{2pt}{2pt}
% ===========================================================

\usepackage{enumitem}
\setlist[itemize]{noitemsep,topsep=2pt,left=0pt .. \parindent}
\pagestyle{empty}
\pdfgentounicode=1

\usepackage{tabularx} % für Tabellen mit automatischer Restbreite (X-Spalte)
\usepackage[hidelinks, bookmarks=false,
            pdftitle={Curriculum Vitae},
            pdfauthor={Name von Arbeit-Suchen}]{hyperref}
\usepackage{xurl}     % Links & E-Mails dürfen an jeder Stelle umbrechen

\newcommand{\cvfont}{\Large\sffamily}

\newlength{\capA}\newlength{\capB}
\settoheight{\capA}{\cvfont H}
\settoheight{\capB}{\normalsize\sffamily\sbseries H}

\begin{document}


\renewcommand{\arraystretch}{1.18}
\noindent\begin{tabularx}{\textwidth}{@{} >{\raggedright\arraybackslash}X >{\sffamily\sbseries}l l @{}}

	\raisebox{\dimexpr\capB-\capA\relax}{\cvfont Curriculum Vitae} & Geburtsdatum: & 10.11.1984 \\
	 & Geburtsort: & Recklinghausen \\
	\smash{\LARGE\sffamily\sbseries Name von Arbeit-Suchen} & Telefon: & 0800 66666 777 \\
	 & Email: & email.a.dresse@mail.com \\
	 & Anschrift: & Irgendeine Straße 77, 07777 Stadt
\end{tabularx}

\section{Ausbildung}
%=================%
\subsection{\href{https://www.beispiel-uni.de}{Irgendeine Universität Berlin} \hfill Okt 2002 -- Sept 2010} % die Dauer des Studiums
M.Sc. Computational Science and Engineering
\hfill Note 2.7 %Abschlussnote

Masterarbeit \flqq \textit{What is \flq Scientific Computing\frq : Its essence and non-essence}\frqq
\hfill Note 1.7

\subsection{\href{https://www.beispiel-uni.de}{Tolle Universität Köln} \hfill Okt 2000 -- Sept 2002}
B.Sc. Mathematik
\hfill Note 3.0

Bachelorarbeit \flqq \textit{Was ist 2+2 und wie wird es bewiesen.}\frqq
\hfill Note 2.7

Auslandssemester im Ausland: \textit{Ausländische Universität}\hfill Sept 2001 -- Feb 2002

\subsection{Schule Recklinghausen \hfill 2002}
Allgemeine Hochschulreife \hfill Note 3.5 %diese Note kann weggelassen werden

\section{Erfahrung}
%=================%
\subsection{\href{https://www.firma.de}{Große Firma} \hfill Stadt}
\subsubsection{Position (60 H. pro Woche) \hfill Sept 2016 -- Dez 2021}
\begin{itemize}
	\item die wichtigste und bedeutendste Sache nennen, die man dort gemacht hat \\
	(\textbf{Cool}, Programm, Sache)
	\item eine andere Sache nennen, die fast genauso bedeutend war
	\item Bereitstellung einer Sache, Lösung eines Problems, Anwendung,\\ Erweiterung etc. Applikation
\end{itemize}

\subsubsection{Position \hfill Jan 2011 -- Aug 2016}
\begin{itemize}
    \item Entwicklung eines Schokoriegels
    \item Betreuung von Praktikanten
\end{itemize}

\subsection{Lehrstuhl für Prokrastination \hfill Irgendeine Universität Berlin}
\subsubsection{Studentische Hilfskraft (10 H. pro Woche) \hfill Jan 2007 -- Aug 2010}
\begin{itemize}
	\item Bücher einscannen, Listen abtippen, Kaffee kochen
	\item Kopieren, Recherchieren, unbezahlte Überstunden machen
\end{itemize}

\section{Fähigkeiten}
%===========================%
\begin{description}[itemsep=1pt, font=\normalfont\sffamily\sbseries]
	\item[Software] \LaTeX, GitHub, MS Office
	\item[Programmiersprachen] Python, Java
	\item[Sprachen] Deutsch (Muttersprache), Englisch (C1), Türkisch (B2)
\end{description}

\end{document}